\scanpage{112}
$H$ et $H'$, on dit qu'ils sont \uline{en position générale} si sur chaque fibre géométrique, les deux paires de points sur la conique isotrope sont disjoints. Mais cela signifie \deleted{justement} qu'il n'y a pas de droite isotrope (sur la dite fibre) orth. \uline{à la fois} à $H$ et $H'$, i.e à $B$. Donc par notre proposition, cela signifie que $\delta_B$ est inv. \deleted{\unsure{i.e.|dire}.} que $\det\big|D$ est inversible. \deleted{\unsure{Si|Le}} Dans ce cas là, \deleted{\unsure{on choisit|choisit}} pour tt choix (possible localmt) \deleted{\unsure{Si on suppose $H,H'$ toutes deux|Si on suppose $H,H'$ étales|Supposons $H,H'$ étales}} de droites iso\unsure{~disjointes|. distinctes}
\[
 D_1,D_2\subset \deleted{\unsure{H}}V\cap\underset{\substack{\shortparallel\\\text{cône isotrope}}}{C},\qquad D'_1,D'_2\subset V'\cap C
\]
on a un rapport anharmonique
\[
 \lambda=\lambda_{(D_1,D_2)}(D'_1,D'_2)\in k
\]
\margin{$\lambda$ \unsure{partout} $\ne 0,1$ i.e $\lambda$ et $1-\lambda$ inv.}
qui est transformé en $1/\lambda$ si on echange $D_1$ et $D_2$ ou $D'_1$ et $D'_2$, et qui en dépend
% archive sheet number: 111
